\chapter{MARCO TEÓRICO}
\label{marco_teorico}

Este capítulo presenta los fundamentos generales empleados en el diseño del
robot móvil diferencial. La exposición avanza desde el robot como sistema
físico (estructura, marcos de referencia y descripción digital) hacia la
plataforma de software que lo integra, los principios de sensado y
actuación, el modelo cinemático y el control de bajo nivel, y finalmente
las capas de mapeo, localización, navegación y exploración autónoma.
\section{El robot móvil: estructura, marcos de referencia y rotaciones}
\label{sec:robot_definicion}

Un robot móvil puede describirse como un conjunto de cuerpos rígidos
(eslabones) unidos mediante articulaciones, organizados como un árbol
cinemático cuya raíz es la base del robot \citep{corke2023}. En una
plataforma diferencial, el árbol se compone del chasis, las dos ruedas
motrices (articulaciones de revolución accionadas) y los elementos
pasivos de apoyo, además de los sensores montados sobre el chasis.
Para describir la posición y orientación del robot se asocia un marco de
referencia a cada cuerpo de interés. La pose de un robot que se desplaza en
el plano queda definida por el vector $q=(x,y,\theta)$, donde $x$ e $y$ son
las coordenadas de la base respecto de un marco fijo (m) y $\theta$ la
orientación (rad). La relación entre un punto expresado en el marco del robot
y el mismo punto en el marco fijo se obtiene mediante la matriz de rotación
en el plano:

\begin{equation}
R(\theta)
=
\begin{bmatrix}
\cos\theta & -\sin\theta\\
\sin\theta & \cos\theta
\end{bmatrix},
\label{eq:rotacion_plana}
\end{equation}

\noindent que, combinada con la traslación $(x,y)$, forma una transformación
homogénea. Encadenando transformaciones se relaciona cualquier par de marcos
del robot; por ejemplo, la pose del sensor láser respecto del marco global se
obtiene componiendo la transformación del mapa a la base del robot con la
transformación fija de la base al sensor \citep{corke2023}. En robótica móvil
es habitual organizar estos marcos en una jerarquía que separa el marco
global del mapa, el marco odométrico (continuo pero con deriva) y el
marco de la base del robot; las convenciones de orientación de ejes (eje $x$
hacia adelante, $y$ hacia la izquierda, $z$ hacia arriba) y de organización
de marcos siguen las recomendaciones REP-103 y REP-105
\citep{rep103,rep105}.

\subsection{Descripción del modelo mediante URDF}
\label{sec:urdf_marco}

El formato URDF (\textit{Unified Robot Description Format}) describe al robot
como un árbol de eslabones y articulaciones. Por cada eslabón se declaran su
geometría visual, su geometría de colisión y sus propiedades inerciales, y por
cada articulación su tipo, sus límites y los marcos que conecta, además de los
sensores y las extensiones de simulación. La misma descripción alimenta la
visualización, la publicación de transformaciones y el simulador, de modo que
el robot real y su modelo digital comparten una única definición geométrica,
parametrizada de forma que dimensiones como el radio de rueda se definan una
sola vez.

El Alg.~\ref{alg:urdf_estructura} resume la organización modular de la
descripción URDF/\texttt{xacro} desarrollada, en la que un archivo raíz
compone el modelo a partir de módulos reutilizables que se seleccionan según
el entorno de ejecución (hardware real o simulación).

\begin{algorithm}[htbp]
\caption{Construcción del modelo URDF/\texttt{xacro} del robot.}
\label{alg:urdf_estructura}
\begin{algorithmic}[1]
\Require \texttt{sim\_mode},\ \texttt{use\_ros2\_control}
\Procedure{DefinirRobot}{} \Comment{\texttt{robot.urdf.xacro}}
\State incluir \texttt{robot\_core.xacro} y \texttt{inertial\_macros.xacro}
\State $\mathcal{L} \gets \{\texttt{base\_link},\ \texttt{chassis},\ \texttt{left\_wheel},\ \texttt{right\_wheel},\ \texttt{caster\_wheel}\}$
\For{cada eslabón $\ell \in \mathcal{L}$}
\State definir geometría visual, colisión e inercia de $\ell$
\EndFor
\For{cada articulación $a$ (fija o continua) entre eslabones de $\mathcal{L}$}
\State conectar eslabón padre e hijo de $a$ y fijar su marco
\EndFor
\State definir \texttt{laser\_frame} y su articulación fija al chasis
\If{\texttt{sim\_mode}}
\State incluir \texttt{lidar\_sim.xacro},\ \texttt{gazebo\_control.xacro}
\Else
\State incluir \texttt{lidar.xacro},\ \texttt{ros2\_control.xacro}
\EndIf
\If{\texttt{use\_ros2\_control}}
\State declarar interfaces de las ruedas: comando (velocidad) y estado (velocidad y posición)
\EndIf
\State \Return modelo del robot
\EndProcedure
\end{algorithmic}
\end{algorithm}

\subsection{Simulación robótica y Gazebo Classic}
\label{sec:simuladores_marco}

Un simulador robótico integra un motor de física, la simulación de sensores y un entorno tridimensional configurable, lo que permite verificar la geometría, la cinemática y la lógica de control antes de operar el hardware. Gazebo Classic, el simulador empleado, se integra de forma nativa con \gls{ROS2}
mediante los paquetes de puente del ecosistema: el modelo URDF se carga en el
mundo simulado y los sensores publican datos con las mismas interfaces que sus
contrapartes físicas \citep{osrf2025gazeboclassic}. La brecha simulación--realidad, es decir fricción, zonas muertas, ruido sensorial y latencias, se considera al interpretar los resultados del Capítulo~\ref{ch:resultados}.

\section{ROS 2: nodos, tópicos y comunicación}
\label{sec:ros2_marco}

\gls{ROS2} es un conjunto de bibliotecas y herramientas de código abierto
para construir sistemas robóticos de forma modular y distribuida
\citep{macenski2022ros2}. Su unidad básica es el \emph{nodo}, un proceso que
encapsula una función específica, y la comunicación principal ocurre mediante
\emph{tópicos}, canales con nombre en los que unos nodos publican mensajes y
otros se suscriben para recibirlos. Cada nodo conoce solo el nombre y el tipo
del tópico, lo que permite reemplazar un componente sin modificar el resto del
sistema. La Fig.~\ref{fig:esquema_nodos} ilustra este patrón.

\begin{figure}[H]
\centering
\setlength{\fboxsep}{6pt}
\shorthandoff{<>}
\fbox{
\begin{tikzpicture}[
>={Stealth[]},
nodo/.style={draw, rounded corners=2pt, minimum width=20mm, minimum height=9mm, align=center, font=\small},
topico/.style={draw, rounded corners=5pt, fill=gray!12, minimum width=18mm, minimum height=9mm, align=center, font=\small},
arista/.style={->, thick},
etiqueta/.style={font=\footnotesize}
]
\node[nodo] (p1) {Nodo\\ publicador};
\node[nodo, below=8mm of p1] (p2) {Nodo\\ publicador};
\node[topico, right=22mm of p1, yshift=-8mm] (t) {Tópico};
\node[nodo, right=22mm of t, yshift=8mm] (s1) {Nodo\\ suscriptor};
\node[nodo, right=22mm of t, yshift=-8mm] (s2) {Nodo\\ suscriptor};
\draw[arista] (p1) -- node[etiqueta, above, sloped]{publica} (t);
\draw[arista] (p2) -- node[etiqueta, below, sloped]{publica} (t);
\draw[arista] (t) -- node[etiqueta, above, sloped]{se suscribe} (s1);
\draw[arista] (t) -- node[etiqueta, below, sloped]{se suscribe} (s2);
\end{tikzpicture}
}
\shorthandon{<>}

\vspace{4pt}

\caption{Patrón de comunicación publicación--suscripción: los nodos publicadores envían mensajes a un tópico y los suscriptores los reciben, de forma anónima y desacoplada.}
\label{fig:esquema_nodos}

\vspace{2pt}

{\footnotesize\textit{Fuente: Elaboración propia.}}
\end{figure}

El transporte se apoya en el estándar DDS (\textit{Data Distribution Service}), que descubre los nodos de la red y negocia políticas de calidad de servicio \citep{macenski2022ros2}. Gracias a ello los nodos pueden repartirse entre varios computadores sin un maestro central, como ocurre aquí entre la Raspberry Pi~4 y la estación de trabajo. \gls{ROS2} ofrece además servicios y acciones, estas últimas usadas para enviar metas de navegación. El árbol de transformaciones mantiene las relaciones espaciales entre los marcos de la Sec.~\ref{sec:robot_definicion}.

\section{Principio de funcionamiento de sensores y actuadores}
\label{sec:sensores_actuadores}

\subsection{Sensor láser LIDAR 2D}
\label{sec:lidar_marco}

Un \gls{LIDAR} mide distancias emitiendo pulsos de luz láser y midiendo la señal reflejada, por tiempo de vuelo o por triangulación óptica \citep{corke2023}. En un \gls{LIDAR} 2D, el emisor y el receptor giran sobre un eje vertical y entregan, en cada revolución, un barrido de pares ángulo y distancia en un plano horizontal, con el que se detectan paredes y obstáculos y se alimentan el SLAM y los mapas de costos \citep{ismail2022eslam}. Sus limitaciones principales son la incapacidad de detectar objetos fuera del plano de barrido y la degradación ante superficies muy absorbentes o reflectantes. El sensor empleado, RPLidar A1, opera por triangulación con alcance nominal de 12~m y cobertura de 360$^\circ$ \citep{SlamtecA1Datasheet}.

\subsection{Encoders incrementales de cuadratura}
\label{sec:encoders_marco}

Un encoder incremental convierte la rotación de un eje en una secuencia de
pulsos eléctricos. En los encoders magnéticos de efecto Hall, un imán
multipolar solidario al eje del motor induce pulsos en dos sensores
desfasados, generando dos canales en cuadratura (desfasados 90$^\circ$). El
orden de los flancos entre ambos canales indica el sentido de giro, y contar
los flancos de ambos canales cuadruplica la resolución básica del sensor
(decodificación~$\times$4). La Fig.~\ref{fig:cuadratura_encoder} muestra ambos
canales, el desfase entre ellos y los flancos contabilizados.

\utemFigurePropia{2.Texto/Imag/fig_cuadratura_encoder.png}{0.85\textwidth}
{Canales en cuadratura de un encoder incremental y decodificación $\times$4.}
{fig:cuadratura_encoder}

Cuando el encoder está montado en el eje del
motor y este acciona la rueda a través de una caja reductora, cada vuelta de
la rueda corresponde a $N_g$ vueltas del motor, por lo que la resolución
teórica en el eje de salida es:

\begin{equation}
N_e^{\mathrm{teo}} = N_{\mathrm{Hall}} \cdot c_q \cdot N_g,
\label{eq:resolucion_encoder}
\end{equation}

\noindent donde $N_e^{\mathrm{teo}}$ es la resolución teórica del encoder
[ticks/vuelta del eje de salida], $N_{\mathrm{Hall}}$ el número de pulsos
por vuelta del eje del motor, $c_q\in\{1,2,4\}$ el factor de decodificación
en cuadratura y $N_g$ la relación de reducción. En la práctica, tolerancias
mecánicas y pérdidas de conteo hacen que la resolución efectiva $N_e$ pueda
diferir del valor teórico, por lo que es habitual determinarla
experimentalmente. Esta resolución es la que relaciona el conteo acumulado
de ticks con la posición y la velocidad angular de cada rueda mediante las
Ecs.~\eqref{eq:posicion_rueda_encoder} y~\eqref{eq:w_ruedas}, base de la
odometría desarrollada en la Sec.~\ref{sec:odometria}; los valores numéricos
del prototipo se consolidan en el Capítulo~\ref{ch:implementacion}.

\subsection{Motores DC y driver de potencia}
\label{sec:motores_marco}

Un motor DC de imanes permanentes gira con velocidad aproximadamente proporcional a la tensión media aplicada y entrega un torque proporcional a la corriente. Como un microcontrolador no puede suministrar esa potencia, se interpone un \emph{driver} basado en un puente H, que aplica la tensión de la batería al motor en ambas polaridades y regula la magnitud de la velocidad mediante \gls{PWM}. Los drivers reales introducen dos efectos relevantes para el control: una caída de tensión interna y una zona muerta a bajos ciclos de trabajo, en la cual el motor no vence la fricción estática y no gira.

\section{Cinemática diferencial y odometría}
\label{sec:odometria}

La odometría estima la pose del robot integrando el desplazamiento medido de
sus ruedas \citep{ferreira2023differential}. Es adecuada para sistemas
embebidos por su bajo costo computacional, pero acumula deriva por errores
de calibración, cuantización de encoders y deslizamiento \citep{corke2023}.

\noindent\textbf{Convención.} $v$ y $\omega$ sin subíndice denotan la
velocidad lineal (m/s) y angular (rad/s) del robot; con subíndice
$j\in\{R,L\}$ denotan las magnitudes correspondientes de la rueda derecha o
izquierda. La posición angular de rueda se denota $\phi_j$ (rad) y el conteo
de encoder se expresa en ticks. El listado completo de símbolos, subíndices
y unidades empleados en el documento se reúne en el \ref{anx:nomenclatura}.

Considérese un robot diferencial con radio de rueda $r$, separación entre
ruedas $b$ y la forma muestreada $q_k=(x_k,y_k,\theta_k)$ de la pose definida
en la Sec.~\ref{sec:robot_definicion}, donde $k$ identifica la muestra
discreta. La relación entre las velocidades angulares de rueda y las
velocidades del robot es la del modelo cinemático diferencial estándar
\citep{corke2023,ferreira2023differential}:

\begin{align}
v[k] &= \frac{r}{2}\left(\omega_R[k]+\omega_L[k]\right),
\label{eq:v_lineal}\\
\omega[k] &= \frac{r}{b}\left(\omega_R[k]-\omega_L[k]\right).
\label{eq:v_angular}
\end{align}

La relación inversa, utilizada para transformar una consigna del robot en
consignas de rueda, es:

\begin{align}
\omega_R^{\mathrm{cmd}}[k]
&=
\frac{v^{\mathrm{cmd}}[k]}{r}
+
\frac{b}{2r}\,\omega^{\mathrm{cmd}}[k],
\notag\\
\omega_L^{\mathrm{cmd}}[k]
&=
\frac{v^{\mathrm{cmd}}[k]}{r}
-
\frac{b}{2r}\,\omega^{\mathrm{cmd}}[k].
\label{eq:cinematica_inversa}
\end{align}

\noindent donde $v^{\mathrm{cmd}}$ y $\omega^{\mathrm{cmd}}$ son las
consignas lineal (m/s) y angular (rad/s) del robot, mientras que
$\omega_R^{\mathrm{cmd}}$ y $\omega_L^{\mathrm{cmd}}$ son las consignas
angulares de las ruedas (rad/s).
Además, una discretización de Euler del modelo cinemático permite propagar la pose
entre muestras:

\begin{align}
x_{k+1} &= x_k + T_k\,v[k]\cos\theta_k,
\label{eq:pose_x}\\
y_{k+1} &= y_k + T_k\,v[k]\sin\theta_k,
\label{eq:pose_y}\\
\theta_{k+1} &= \theta_k + T_k\,\omega[k].
\label{eq:pose_theta}
\end{align}

\noindent donde $T_k$ es el intervalo temporal entre las muestras $k$ y
$k+1$ (s). Estas ecuaciones expresan la integración odométrica del modelo
cinemático diferencial general \citep{ferreira2023differential}. El integrador odométrico efectivamente ejecutado pertenece a la
capa de control descrita en el Capítulo~\ref{ch:implementacion}.

\subsection{Estimación de estados de rueda a partir de encoders}
\label{sec:estados_rueda}

A partir del conteo acumulado de encoder $c_j[k]$ (ticks), la posición
angular de cada rueda se obtiene mediante:

\begin{equation}
\phi_j[k]
=
\frac{2\pi}{N_e}\,c_j[k],
\qquad j\in\{R,L\},
\label{eq:posicion_rueda_encoder}
\end{equation}

\noindent donde $N_e$ es la resolución efectiva del encoder (ticks/vuelta),
definida en la Sec.~\ref{sec:encoders_marco}. La velocidad angular se estima por diferencias finitas entre dos lecturas
consecutivas:

\begin{equation}
\widehat{\omega}_j[k]
=
\frac{\phi_j[k]-\phi_j[k-1]}{\Delta t_k},
\qquad j\in\{R,L\}.
\label{eq:w_ruedas}
\end{equation}

\noindent donde $\widehat{\omega}_j[k]$ es la velocidad angular estimada de
la rueda $j$ (rad/s) y $\Delta t_k$ el tiempo real transcurrido entre las
lecturas empleadas en segundos. La definición por casos de $\Delta t_k$ y la
secuencia de lectura ejecutada por la interfaz de hardware se desarrollan en
el Capítulo~\ref{ch:implementacion} y en el
\ref{anx:implementacion_diffdrive}.

\section{Identificación FOPDT}
\label{sec:fodt}

Para obtener un modelo destinado a la sintonización del controlador se
identifica, para cada rueda, una estructura de primer orden con retardo FOPDT
(\textit{First Order Plus Dead Time}) a partir de la respuesta experimental
ante un escalón de \gls{PWM} \citep{maxim2023fopdt}:

\begin{equation}
G_j(S)
=
\frac{K_{m,j}}{1+\tau_j S}\,
\mathrm{e}^{-L_{d,j}S},
\qquad j\in\{R,L\}.
\label{eq:fodt}
\end{equation}

\noindent donde $G_j(S)$ es la función de transferencia identificada para la
rueda $j$, $K_{m,j}$ la ganancia estática [(m/s)/PWM], $\tau_j$ la constante
de tiempo (s) y $L_{d,j}$ el retardo puro o tiempo muerto (s). La variable
compleja de Laplace se denota con mayúscula, $S$, para distinguirla a simple
vista del segundo (s), unidad de tiempo escrita en redonda; los productos
$\tau_j S$ y $L_{d,j}\,S$ son adimensionales.
Los parámetros $(K_m,\tau,L_d)$ se estiman minimizando el error
cuadrático entre la respuesta medida y la respuesta del modelo. La
discretización, la cuantización del retardo y las aproximaciones numéricas
empleadas por los scripts se documentan en el
\ref{anx:fodt_scripts}.

\section{Control PID}
\label{sec:pid}

El control de velocidad de cada rueda se organiza como un lazo cerrado: la
velocidad medida por el encoder se compara con la referencia, y la diferencia
alimenta al controlador, que ajusta la tensión media entregada al motor. La
Fig.~\ref{fig:lazo_pid} resume esta estructura.

\utemFigurePropia{2.Texto/Imag/fig_lazo_pid.png}{0.85\textwidth}
{Lazo de control de velocidad de una rueda con controlador PID.}
{fig:lazo_pid}

El controlador de velocidad se ejecuta de forma independiente para cada rueda en el Arduino Nano. La forma continua del PID, tomada de la formulación clásica de control realimentado \citep{corke2023}, solo se utiliza como referencia conceptual:

\begin{equation}
u_j(t)
=
K_{p,j}^{\mathrm{cont}}e_j(t)
+
K_{i,j}^{\mathrm{cont}}\int_0^t e_j(\sigma)\,d\sigma
+
K_{d,j}^{\mathrm{cont}}\frac{de_j(t)}{dt}.
\label{eq:pid_continuo}
\end{equation}

\noindent donde $t$ es el tiempo continuo, $u_j(t)$ la acción de control, $e_j(t)$ el error de velocidad, $K_{p,j}^{\mathrm{cont}}$, $K_{i,j}^{\mathrm{cont}}$ y $K_{d,j}^{\mathrm{cont}}$ las ganancias proporcional, integral y derivativa de la forma continua, respectivamente, y $\sigma$ la variable de integración.

El controlador no se evalúa de forma continua, sino a intervalos regulares en
el microcontrolador. Su formulación discreta específica, la realización en el
firmware y las ganancias definitivas se desarrollan en la
Sec.~\ref{sec:motorcontrol_pid} y en el \ref{anx:pid_ecuaciones}.

\section{Sintonización del controlador PID mediante PSO}
\label{sec:sintonizacion}

Las ganancias del PID se obtienen mediante un procedimiento híbrido:
identificación FOPDT por rueda (Sec.~\ref{sec:fodt}), búsqueda de ganancias
iniciales mediante optimización por enjambre de partículas y refinamiento
experimental posterior sobre el hardware. Como referencia analítica
complementaria se emplea además la sintonía Lambda, cuyo desarrollo y
condiciones de comparación se documentan en el \ref{anx:lambda_detalle}.

\subsection{Optimización por enjambre de partículas}
\label{sec:pso_marco}

PSO (\textit{Particle Swarm Optimization}) es una metaheurística
poblacional inspirada en el comportamiento colectivo de enjambres. En este
trabajo se emplea la variante con factor de constricción, en la formulación
aplicada a la sintonización de controladores PID en robótica móvil por
\citet{gokcce2021particle}. Cada partícula $i$ representa un candidato de solución, en este caso, un
conjunto de ganancias PID, mediante una posición $\xi_{i,d}$ y una
velocidad $V_{i,d}$, ambas adimensionales, con $d\in\{1,2,3\}$ asociado a
las tres ganancias. En cada generación $g$, la
velocidad de cada partícula se actualiza combinando su inercia, la atracción
hacia su mejor posición personal y la atracción hacia la mejor posición
global del enjambre:

\begin{equation}
V_{i,d}^{g+1}
=
\zeta
\left[
w_{\mathrm{in}}^{g}V_{i,d}^{g}
+c_1\rho_{1,i,d}^{g}
\left(P_{i,d}-\xi_{i,d}^{g}\right)
+c_2\rho_{2,i,d}^{g}
\left(G_d^{g}-\xi_{i,d}^{g}\right)
\right],
\label{eq:pso_vel_marco_teorico}
\end{equation}

\noindent donde $P_{i,d}$ es la mejor posición personal almacenada,
$G_d^{g}$ la componente del mejor global, $w_{\mathrm{in}}^{g}$ el peso de
inercia, $c_1$ y $c_2$ los coeficientes cognitivo y social,
$\rho_{1,i,d}^{g}$ y $\rho_{2,i,d}^{g}$ números aleatorios uniformes en
$[0,1]$ y $\zeta$ el factor de constricción, que garantiza la convergencia
del enjambre. Los valores utilizados para $c_1$, $c_2$ y $\zeta$ se consolidan
en la Tabla~\ref{tab:pso_config}. El peso de inercia decrece linealmente con
las generaciones:

\begin{equation}
w_{\mathrm{in}}^{g}
=
\frac{g_{\max}-g}{g_{\max}},
\qquad
g=1,\ldots,g_{\max},
\label{eq:pso_inercia}
\end{equation}

\noindent donde $g_{\max}$ es el número máximo de generaciones; de esta
forma el enjambre privilegia la exploración al inicio y la explotación al
final. La función de costo empleada para la sintonización combina el error
cuadrático medio de la velocidad respecto de la referencia con una
penalización por sobreimpulso, favoreciendo respuestas rápidas pero
amortiguadas. Las escalas de búsqueda, las cotas, la configuración del
enjambre y las particularidades numéricas de la ejecución se presentan en el
Capítulo~\ref{ch:implementacion} y en el \ref{anx:pso_detalle}.

\section{SLAM basado en grafos}
\label{sec:slam}

La odometría pura acumula deriva, por lo que se requiere SLAM para estimar
simultáneamente la trayectoria del robot y el mapa del entorno
\citep{thrun2005probabilistic,ngoc2023efficient}. En el enfoque basado en
grafos, cada nodo del grafo representa una pose del robot y cada arista una
restricción relativa entre dos poses, obtenida de la odometría o del ajuste
de barridos láser (\textit{scan matching}), que alinea geométricamente dos
barridos consecutivos para estimar el desplazamiento entre ellos. Cuando el
robot vuelve a observar una zona ya visitada se agrega una restricción de
\emph{cierre de lazo}, y la optimización del grafo, un problema de mínimos
cuadrados no lineales, redistribuye el error acumulado a lo largo de toda
la trayectoria, corrigiendo el mapa y la transformación entre el marco del
mapa y el marco odométrico. La Fig.~\ref{fig:grafo_poses} representa esta
estructura.

\utemFigurePropia{2.Texto/Imag/fig_grafo_poses.png}{0.80\textwidth}
{Grafo de poses: nodos, restricciones relativas y cierre de lazo.}
{fig:grafo_poses}

El procesamiento asíncrono de los barridos, la detección de cierre de lazo y
un consumo compatible con computadores embebidos son las características que
hacen viable esta técnica en plataformas de bajo costo. La técnica utilizada
en este trabajo es SLAM Toolbox, una realización del SLAM basado en grafos de
poses que construye el mapa de ocupación de forma incremental durante la
operación, mantiene y optimiza el grafo mientras el robot se desplaza, y
permite almacenar el mapa resultante para reutilizarlo después en la
localización \citep{macenski2021slamtoolbox}. La comparación con las demás
alternativas de SLAM 2D se presenta en el Capítulo~\ref{ch:estado_arte}, y la
configuración concreta empleada en la Sec.~\ref{sec:slam_ch9}.

\section{Localización en mapas conocidos}
\label{sec:localizacion_marco}

Una vez construido y almacenado el mapa del entorno, la estimación de la pose
puede realizarse sin volver a ejecutar el proceso de mapeo. Para este propósito
se emplea la localización de Monte Carlo adaptativa, un método de localización
probabilística basado en un filtro de partículas
\citep{thrun2005probabilistic}.
Cada partícula representa una hipótesis de pose. Las hipótesis se propagan con la odometría y se ponderan según la concordancia entre los barridos del \gls{LIDAR} y el mapa de ocupación almacenado, y luego el remuestreo concentra las partículas en las regiones de mayor probabilidad. Este método estima la transformación entre el marco del mapa y el odométrico, mientras que la odometría mantiene la transformación entre este último y la base del robot. A diferencia del SLAM de la Sec.~\ref{sec:slam}, la localización no construye ni modifica el mapa, sino que sitúa al robot dentro de uno ya conocido. El nodo utilizado y su configuración se detallan en la Sec.~\ref{sec:amcl_ch9}.

\section{Planificación de rutas y control local}
\label{sec:nav_marco}

La navegación autónoma opera sobre una representación del entorno en forma de
rejilla de ocupación, donde cada celda clasifica el espacio como libre, ocupado
o desconocido \citep{lavalle2006planning}. Sobre esa rejilla se construyen mapas de costo que agregan a los obstáculos una zona de inflación que penaliza
transitar cerca de ellos \citep{nav2_doc}. Sobre esta representación, el planificador global calcula una ruta desde la posición actual hasta la
meta. En este trabajo la búsqueda global se resuelve mediante el algoritmo A$^*$ sobre la rejilla derivada del mapa de costos. El planificador empleado y la opción de configuración que habilita esta búsqueda se presentan en el
Capítulo~\ref{ch:implementacion}. A$^*$ ordena la expansión de
los nodos mediante la función de evaluación:

\begin{equation}
f(n)=g(n)+h(n),
\label{eq:astar}
\end{equation}

\noindent donde $n$ es el nodo evaluado, $g(n)$ representa el costo acumulado
desde el origen hasta dicho nodo y $h(n)$ una estimación del costo restante
hasta la meta. La implementación de Nav2 incorpora además los costos de su propia representación, cuyo escalamiento interno no se reconstruyó en este trabajo. Por eso no se emplean la admisibilidad de la heurística ni la optimalidad de la ruta como criterio de validación, sino la capacidad de generar una ruta global que el controlador local pueda seguir.

El control local convierte la ruta global en comandos de velocidad
ejecutables. El enfoque de ventana dinámica considera únicamente las
velocidades $(v,\omega)$ alcanzables desde el estado actual dado el límite
de aceleración del robot, simula hacia adelante las trayectorias asociadas a
cada candidato y las evalúa mediante una suma ponderada de criterios:

\begin{equation}
J_{\mathrm{DWB}}(v,\omega)
=
\sum_{c\in\mathcal{C}}
S_c\,C_c(v,\omega),
\label{eq:costo_ctrl_local}
\end{equation}

\noindent donde $J_{\mathrm{DWB}}(v,\omega)$ es el puntaje total de una
trayectoria candidata con velocidades lineal $v$ y angular $\omega$;
$\mathcal{C}$ es el conjunto de criterios de evaluación, tales como
proximidad a obstáculos, alineación con la ruta y avance hacia la meta,
$C_c(v,\omega)$ el costo asignado por el criterio $c$ y $S_c$ su peso. Se
selecciona el candidato de menor puntaje, lo que permite seguir la ruta
global mientras se reacciona ante obstáculos cercanos. La navegación de alto
nivel se orquesta mediante un árbol de comportamiento: una estructura
modular que encadena la planificación y el seguimiento, y que ante fallos
activa acciones de recuperación de complejidad creciente sin modificar la
lógica del resto del sistema. Los componentes concretos, sus criterios
configurados y las acciones de recuperación se presentan en el
Capítulo~\ref{ch:implementacion}.

\section{Exploración autónoma basada en fronteras}
\label{sec:exploracion_marco}

La exploración autónoma permite que el robot seleccione sus propias metas
mientras construye el mapa. En el enfoque basado en fronteras, introducido por \citet{yamauchi1997frontier}
y todavía vigente en formulaciones recientes \citep{liu2025frontier}, una
\emph{celda de frontera} es una celda desconocida adyacente  a espacio libre conocido; las celdas de frontera
contiguas se agrupan en fronteras, y cada frontera se caracteriza por un
punto representativo, su tamaño y su distancia al robot. Las fronteras se
ordenan mediante un puntaje que combina la distancia, privilegiando metas
cercanas, y el tamaño, privilegiando zonas que aportan más información,
y la mejor frontera se envía como meta al sistema de navegación. Cuando una
meta no muestra progreso dentro de un tiempo límite, se incorpora a una
lista negra para evitar reintentos improductivos, y el proceso se repite
hasta que no quedan fronteras alcanzables, momento en que el mapa se
considera completo. Los parámetros y particularidades de la implementación
utilizada se presentan en el Capítulo~\ref{ch:implementacion} y en el
\ref{anx:explore_impl}.

\section*{Conclusión del capítulo}

El capítulo estableció, en un nivel general, los fundamentos de cada elemento
del sistema enumerados en su introducción, desde la descripción física del
robot hasta la exploración por fronteras. Sobre estos fundamentos, el
Capítulo~\ref{ch:implementacion} y los anexos desarrollan los algoritmos,
parámetros y configuraciones con que cada elemento se materializa en el
prototipo.